\section{Monte Carlo and Data Samples}
\label{sec:llvv_samples}

The specific simulated samples and empirical datasets utilized throughout this analysis are outlined in this section. Captured by the ATLAS detector between 2015 and 2018, the Run~2 $pp$ collision dataset operates at a center-of-mass energy of $\sqrt{s} = 13~\text{TeV}$.

By merging detailed detector response models with theoretical predictions, Monte Carlo (MC) simulations serve a critical function in profiling the topological and kinematic traits of the signal, while also quantifying the expected yields from all major SM backgrounds.

Maintaining consistency in event selection rules, calibration routines, and physics object reconstruction across both simulated samples and actual data remains a cornerstone of this study. Such a unified methodology guarantees an impartial, direct alignment between theoretical models and observed events, ultimately securing the reliability of the extracted physics outcomes.

The subsequent subsections initially present an exhaustive summary of the MC models generated for various background and signal processes. Following this, the discussion breaks down the exact data samples, the applied data quality restrictions, and the specific trigger logic leveraged to isolate analysis events. Exhaustive inventories containing dataset identifiers (DSIDs), trigger chains, generator parameters, and Good Run Lists are cataloged in Appendix~\ref{subsec:samples_zz}.

\subsection{The ATLAS Simulation Framework}
Generating simulated events within the ATLAS collaboration adheres to a highly standardized sequential pipeline, which reliably mirrors proton--proton interactions and the subsequent hardware responses. This simulation workflow is anchored by three primary phases:

\begin{enumerate}
    \item \textbf{Event Generation:} The fundamental hard-scatter collision is replicated utilizing an array of event generators. These specialized programs calculate the matrix element (ME) governing a given physical interaction, typically achieving leading order (LO) or next-to-leading order (NLO) accuracy within perturbative QCD. Subsequently, this ME evaluation is linked to a parton shower (PS) framework—such as the integrated module within \textsc{Sherpa}~\cite{Bothmann_2019} or the standalone \textsc{Pythia}~8~\cite{Sjostrand:2014zea} algorithm. This PS stage simulates particle decays, hadronization, multiple parton interactions, as well as final- and initial-state radiation. For every individual sample, the designated PDF (Parton Distribution Function) set, perturbative precision, and specific generator are carefully selected to yield the highest-fidelity process modeling.

    \item \textbf{Detector Simulation:} Utilizing the \textsc{Geant4} software suite~\cite{GEANT4:2002zbu}, the stable particles produced during the initial generation phase are tracked as they traverse a meticulous digital replica of the ATLAS detector's material makeup and spatial geometry. This crucial step models the exact physical interactions between the traversing particles and the detector's passive and active components, generating simulated ``hits'' and energy deposits across the sub-systems.

    \item \textbf{Digitisation and Reconstruction:} These artificial detector hits are subsequently translated into digitized electronic pulses, closely mimicking authentic hardware readouts. To accurately reflect the LHC's intense high-luminosity conditions, numerous inelastic $pp$ collisions—commonly referred to as pile-up—are superimposed onto every simulated hard-scatter event. To model these pile-up interactions, \textsc{Pythia}~8 is employed alongside the \textsc{NNPDF2.3LO} PDF set~\cite{Ball:2012cx} and the A3 tune~\cite{ATL-PHYS-PUB-2016-017}. Furthermore, the simulated pile-up multiplicity distribution is explicitly reweighted to mirror the exact profile observed in the empirical data. As a final step, the simulated digital output is fed through the exact same reconstruction software used for genuine collision data, ultimately assembling high-level physics objects like jets, missing transverse momentum ($E_{\text{T}}^{\text{miss}}$), electrons, and muons.
\end{enumerate}

To reconcile minor discrepancies in identification, reconstruction, and trigger efficiencies between the MC simulations and empirical data, specific scale factors (or correction weights) are systematically applied to all simulated datasets.


\subsection{Monte Carlo Samples}

Deriving accurate background estimations and signal efficiencies necessitates a broad portfolio of MC event samples. These datasets are explicitly generated to align with the operating conditions of the 2015--2018 data-taking campaigns, and they undergo the exact same calibration algorithms and reconstruction software pipelines as the physical collision data. This rigid symmetry ensures that all physics objects are handled consistently, authorizing a seamless comparison between theoretical expectations and experimental observations. The ensuing segments outline the specialized generators tasked with simulating both background and signal activities. 


\subsubsection{Signal Samples}
This study is fundamentally geared toward measuring the inclusive SM $\text{ZZ}$ production cross-section within the $\ell^+\ell^-\nu\bar{\nu}$ decay topology, simultaneously probing for BSM phenomena via anomalous gauge boson couplings.

\paragraph{Standard Model $\text{ZZ}$ Production}
The dominant signal topology involves the creation of paired $\text{Z}$ bosons that subsequently decay into the $\ell^+\ell^-\nu\bar{\nu}$ final state. This encompasses loop-mediated gluon-gluon fusion as well as quark-antiquark annihilation pathways, both of which are modeled utilizing \textsc{Sherpa}~2.2.2. Conversely, the electroweak generation of $\text{ZZ}$ accompanied by two jets is simulated independently using \textsc{MadGraph5}~\cite{Alwall:2014hca}. 
Table~\ref{tab:sherpa-xs} displays the resulting MC-derived cross-sections for both the $\text{ZZ}jj$ and $\text{ZZ}$ fiducial phase spaces. An exhaustive inventory of these generated samples—featuring their associated k-factors and DSIDs—can be referenced in Appendix~\ref{subsec:samples_zz} (specifically Table~\ref{tab:ZZ_signal_samples}).

\begin{table}[htbp]
\centering
\small
\renewcommand{\arraystretch}{1.4}
\begin{tabular}{l c}
\hline
\textbf{Cross-section}  & \textbf{Calculated value [fb] } \\
\hline
$\sigma^{\text{fid}}_{\text{ZZ} \to \ell\ell\nu\nu}$ & $20.08 \pm 1.0_{\text{scale}} \pm 0.30_{\text{PDF}} \pm 0.24_{\alpha_S}$ \\
$\sigma^{\text{fid}}_{\text{ZZ}jj \to \ell\ell\nu\nu jj}$ & $0.94 \pm 0.19_{\text{scale}} \pm 0.01_{\text{PDF}} \pm 0.02_{\alpha_S}$ \\
\hline
\end{tabular}
\caption{Fiducial cross-sections for $\text{ZZ} \to 2\ell 2\nu$ and $\text{ZZ}jj \to 2\ell 2\nu jj$ productions at $\sqrt{s} = 13~\text{TeV}$. The $\text{ZZ}$ and $\text{ZZ}jj$ cross-sections are obtained with the QCD production simulated using \textsc{Sherpa} and the electroweak production simulated using \textsc{MadGraph5} (NLO EW). Theoretical uncertainties from QCD scale variation, PDFs, and $\alpha_S$ are quoted.}
\label{tab:sherpa-xs}
\end{table}

\paragraph{Samples with Anomalous Gauge Couplings}
The hunt for new physics relies on evaluating anomalous quartic gauge couplings (aQGC) alongside anomalous neutral triple gauge couplings (nTGC) through the lens of an Effective Field Theory (EFT) framework. To accurately replicate the kinematic distortions introduced by these BSM operators, a specialized batch of MC samples was synthesized. Comprehensive DSID registries cataloging the aQGC and nTGC search models are provided in Tables~\ref{tab:ZZEFTsamples} and~\ref{tab:ZZaQGCsamples} within Appendix~\ref{subsec:samples_zz}.


\subsubsection{Background Samples}
Contaminating backgrounds infiltrating the $\ell^+\ell^- + E_{\text{T}}^{\text{miss}}$ target state are generally classified into two camps: reducible backgrounds triggered by severe detector mismeasurements or misidentified objects, and irreducible backgrounds characterized by genuine $E_{\text{T}}^{\text{miss}}$ paired with prompt leptons. Exhaustive data tables cataloging these background samples are available in Appendix~\ref{subsec:samples_zz}.

\paragraph{Diboson and Triboson Backgrounds}
The most prominent irreducible background interference stems from the generation of $\text{WW}$, $\text{WZ}$, and fully leptonic $\text{ZZ} \to \ell^+\ell^-\ell^+\ell^-$ events. Minor contamination originating from triboson processes (such as $\text{ZZZ}$, $\text{WZZ}$, $\text{WWZ}$, and $\text{WWW}$) is similarly accounted for. These specific interactions are primarily simulated via \textsc{Powheg}+\textsc{Pythia}~8 and \textsc{Sherpa} (refer to Tables~\ref{tab:ZZbkg_samples},~\ref{tab:WZandWWsamples}, and~\ref{tab:VVVsamples}).

\paragraph{Top Quark Backgrounds}
Another dominant background source arises from single-top and top-antitop ($t\bar{t}$) generation, specifically leaking through their dileptonic decay branches. Top quark production occurring in association with a vector boson ($t\bar{t}\text{V}$) provides an additional contribution. To model these events, \textsc{MadGraph5\_aMC@NLO} and \textsc{Powheg}+\textsc{Pythia}~8 are deployed (see Tables~\ref{tab:Topsamples} and~\ref{tab:ttVsamples}).

\paragraph{$\text{Z}/\gamma^*$+jets Background}
Acting as the primary reducible background, this contamination originates from Drell-Yan processes plagued by artificially inflated, mismeasured $E_{\text{T}}^{\text{miss}}$. The \textsc{Sherpa}~2.2.1 generator, equipped with NLO-accurate matrix elements, is tasked with modeling this behavior. These specific samples are partitioned by lepton flavor and generated across discrete slices of the partonic transverse momentum sum (documented in Tables~\ref{tab:ZjetsSamplesSherpa1},~\ref{tab:ZjetsSamplesSherpa2}, and~\ref{tab:ZjetsSamplesSherpa3}).

\paragraph{Higgs Boson Backgrounds}
Interactions resulting in the creation of a Higgs boson can occasionally mimic the targeted final state. These relatively minor background contributions are replicated utilizing \textsc{Powheg}+\textsc{Pythia}~8 (detailed in Table~\ref{tab:Higgssamples}).


\subsection{Data Samples}

The empirical dataset is analyzed utilizing the conventional ATLAS Release 21 software environment. Stringent filters are applied so that only events captured while all detector sub-systems were fully operational and beam conditions remained perfectly stable are admitted. To enforce this, the official Run~2 Standard Good Run Lists (GRL), curated by the ATLAS Data Quality team for the 2015--2018 operational windows, are utilized, securing an integrated luminosity of $140~\text{fb}^{-1}$. The precise GRL file directories referenced for each distinct data-taking epoch are recorded in Table~\ref{tab:grls_app} of Appendix~\ref{subsec:samples_zz}.

To guarantee a highly resilient and encompassing acceptance rate for events featuring at least one energetic lepton, the selection algorithm enforces a logical disjunction (OR) spanning multiple trigger paths. By fusing di-lepton-seeded triggers with single-lepton variants, this comprehensive trigger strategy maximizes event recovery across an expansive kinematic phase space. 

Implementing this dual-trigger approach is essential to salvage signal events where the transverse momentum of an individual lepton might dip below the activation threshold required by the dominant high-$p_{\text{T}}$ single-lepton triggers. Consequently, events are harvested using a fused array of un-prescaled single-muon and single-electron triggers. Depending on the applied isolation and identification strictness (e.g., tight or medium likelihood criteria), lepton flavor, and the specific data-taking year, the lowest transverse momentum floors for these triggers fluctuate between 20~GeV and 26~GeV. For leptons exhibiting substantially higher transverse momenta (for instance, exceeding 140~GeV or 50~GeV), supplementary triggers featuring relaxed identification parameters are activated to further patch efficiency gaps. A complete and exhaustive breakdown of the specific software trigger chains deployed throughout this study is laid out in Table~\ref{tab:lepton_triggers_app} of Appendix~\ref{subsec:samples_zz}.